% ============================================================
%  MTH636 – Game Theory | Lecture 21
%  Kuhn's Theorem, SPE Existence & Applications
% ============================================================
\documentclass[12pt, a4paper]{article}

% ============================================================
%  MTH636 -- Game Theory
%  Notes Preamble — Color-Coded, TikZ-Enhanced
% ============================================================

% --- Core packages ---
\usepackage[T1]{fontenc}
\usepackage[utf8]{inputenc}
\usepackage{lmodern}
\usepackage[margin=2.5cm, top=2.8cm, bottom=2.8cm]{geometry}
\usepackage{microtype}

% --- Math ---
\usepackage{amsmath, amssymb, amsthm, mathtools}
\usepackage{bm}        % bold math
\usepackage{cancel}    % strike-through in math

% --- Color & Boxes ---
\usepackage[dvipsnames, svgnames, x11names, table]{xcolor}
\usepackage[most]{tcolorbox}
\tcbuselibrary{breakable, skins, theorems}

% --- TikZ & Diagrams ---
\usepackage{tikz}
\usepackage{tikz-qtree}
\usetikzlibrary{arrows.meta, positioning, shapes.geometric, fit, backgrounds,
                decorations.pathreplacing, calc, matrix, trees}
\usepackage{pgfplots}
\pgfplotsset{compat=1.18}

% --- Tables ---
\usepackage{booktabs}
\usepackage{multirow}
\usepackage{array}
\usepackage{tabularx}
\usepackage{makecell}

% --- Typography & Layout ---
\usepackage{enumitem}
\usepackage{parskip}
\usepackage{fancyhdr}
\usepackage{titlesec}
\usepackage{titletoc}
\usepackage{hyperref}
\usepackage{graphicx}
\usepackage{float}
\usepackage{caption}

% --- Color Palette ---
\definecolor{defblue}{HTML}{1A5276}      % Definition header
\definecolor{defbluebg}{HTML}{D6EAF8}    % Definition background
\definecolor{thmgreen}{HTML}{1E8449}     % Theorem header
\definecolor{thmbg}{HTML}{D5F5E3}        % Theorem background
\definecolor{exorange}{HTML}{B7770D}     % Example header
\definecolor{exbg}{HTML}{FEF9E7}         % Example background
\definecolor{warnred}{HTML}{922B21}      % Warning/Key Insight header
\definecolor{warnbg}{HTML}{FADBD8}       % Warning background
\definecolor{noticegray}{HTML}{2C3E50}   % Note header
\definecolor{noticebg}{HTML}{EBF5FB}     % Note background
\definecolor{taketawayyellow}{HTML}{7D6608} % Key Takeaway header
\definecolor{takeawaybg}{HTML}{FDFEFE}   % Takeaway background
\definecolor{accentpurple}{HTML}{6C3483} % Accent / proof
\definecolor{lightgray}{HTML}{F2F3F4}    % shading

% ============================================================
%  TCOLORBOX ENVIRONMENTS
% ============================================================

% -- Definition (Blue) --
\newtcolorbox{defbox}[2][]{
  enhanced, breakable,
  colback=defbluebg, colframe=defblue, coltitle=white,
  fonttitle=\bfseries\sffamily,
  title={Definition: #2},
  attach boxed title to top left={yshift=-2mm, xshift=4mm},
  boxed title style={colback=defblue, rounded corners},
  arc=4pt, boxrule=1.2pt,
  left=6pt, right=6pt, top=8pt, bottom=6pt,
  #1
}

% -- Theorem / Proposition (Green) --
\newtcolorbox{thmbox}[2][]{
  enhanced, breakable,
  colback=thmbg, colframe=thmgreen, coltitle=white,
  fonttitle=\bfseries\sffamily,
  title={Theorem: #2},
  attach boxed title to top left={yshift=-2mm, xshift=4mm},
  boxed title style={colback=thmgreen, rounded corners},
  arc=4pt, boxrule=1.2pt,
  left=6pt, right=6pt, top=8pt, bottom=6pt,
  #1
}

% -- Proposition (Green variant) --
\newtcolorbox{propbox}[2][]{
  enhanced, breakable,
  colback=thmbg, colframe=thmgreen!70!black, coltitle=white,
  fonttitle=\bfseries\sffamily,
  title={Proposition: #2},
  attach boxed title to top left={yshift=-2mm, xshift=4mm},
  boxed title style={colback=thmgreen!70!black, rounded corners},
  arc=4pt, boxrule=1.2pt,
  left=6pt, right=6pt, top=8pt, bottom=6pt,
  #1
}

% -- Example (Orange/Yellow) --
\newtcolorbox{exbox}[2][]{
  enhanced, breakable,
  colback=exbg, colframe=exorange, coltitle=white,
  fonttitle=\bfseries\sffamily,
  title={Example: #2},
  attach boxed title to top left={yshift=-2mm, xshift=4mm},
  boxed title style={colback=exorange, rounded corners},
  arc=4pt, boxrule=1.2pt,
  left=6pt, right=6pt, top=8pt, bottom=6pt,
  #1
}

% -- Key Insight / Warning (Red) --
\newtcolorbox{warnbox}[2][]{
  enhanced, breakable,
  colback=warnbg, colframe=warnred, coltitle=white,
  fonttitle=\bfseries\sffamily,
  title={Key Insight: #2},
  attach boxed title to top left={yshift=-2mm, xshift=4mm},
  boxed title style={colback=warnred, rounded corners},
  arc=4pt, boxrule=1.2pt,
  left=6pt, right=6pt, top=8pt, bottom=6pt,
  #1
}

% -- Note (Gray-Blue) --
\newtcolorbox{notebox}[1][]{
  enhanced, breakable,
  colback=noticebg, colframe=noticegray!60, coltitle=white,
  fonttitle=\bfseries\sffamily,
  title={Note},
  attach boxed title to top left={yshift=-2mm, xshift=4mm},
  boxed title style={colback=noticegray!60, rounded corners},
  arc=4pt, boxrule=0.8pt,
  left=6pt, right=6pt, top=8pt, bottom=6pt,
  #1
}

% -- Key Takeaway --
\newtcolorbox{takebox}[1][]{
  enhanced, breakable,
  colback=takeawaybg, colframe=taketawayyellow, coltitle=white,
  fonttitle=\bfseries\sffamily,
  title={Key Takeaway},
  attach boxed title to top left={yshift=-2mm, xshift=4mm},
  boxed title style={colback=taketawayyellow, rounded corners},
  arc=4pt, boxrule=1.2pt,
  left=6pt, right=6pt, top=8pt, bottom=6pt,
  #1
}

% -- Intuition --
\newtcolorbox{intubox}[1][]{
  enhanced, breakable,
  colback=accentpurple!8, colframe=accentpurple!60, coltitle=white,
  fonttitle=\bfseries\sffamily,
  title={Intuition},
  attach boxed title to top left={yshift=-2mm, xshift=4mm},
  boxed title style={colback=accentpurple!80, rounded corners},
  arc=4pt, boxrule=0.8pt,
  left=6pt, right=6pt, top=8pt, bottom=6pt,
  #1
}

% ============================================================
%  PAYOFF MATRIX HELPER
%  Usage: inside tabular, use \payoff{a}{b} for (a,b)
% ============================================================
\newcommand{\payoff}[2]{$(#1,\;#2)$}
\newcommand{\nash}[2]{\underline{$#1$},\underline{$#2$}}  % underline NE payoffs
\newcommand{\Nset}{\mathbb{N}}
\newcommand{\Rset}{\mathbb{R}}
\newcommand{\Sset}{\mathcal{S}}
\newcommand{\Ui}{u_i}

% ============================================================
%  SECTION FORMATTING
% ============================================================
\titleformat{\section}
  {\large\bfseries\sffamily\color{defblue}}
  {\thesection}{1em}{}[\titlerule]

\titleformat{\subsection}
  {\normalsize\bfseries\sffamily\color{thmgreen!80!black}}
  {\thesubsection}{1em}{}

\titleformat{\subsubsection}
  {\normalsize\itshape\sffamily\color{exorange!80!black}}
  {\thesubsubsection}{1em}{}

% ============================================================
%  HEADER / FOOTER
% ============================================================
\pagestyle{fancy}
\fancyhf{}
\fancyhead[L]{\sffamily\small\color{defblue} MTH636 -- Game Theory}
\fancyhead[R]{\sffamily\small\color{defblue} \leftmark}
\fancyfoot[C]{\sffamily\small\thepage}
\renewcommand{\headrulewidth}{0.4pt}

% ============================================================
%  HYPERREF SETUP
% ============================================================
\hypersetup{
  colorlinks=true,
  linkcolor=defblue,
  urlcolor=thmgreen,
  citecolor=accentpurple,
  bookmarksnumbered=true,
  pdftitle={MTH636 Game Theory Notes},
  pdfauthor={MTH636 Student Notes}
}

% ============================================================
%  MISC HELPERS
% ============================================================
\newcommand{\BR}{\text{BR}}        % Best Response
\newcommand{\NE}{\text{NE}}        % Nash Equilibrium
\newcommand{\SPNE}{\text{SPNE}}    % Subgame Perfect NE
\newcommand{\Dom}{\text{dom}}      % Dominance
\newcommand{\mixs}[1]{\sigma_{#1}} % Mixed strategy
\newcommand{\pures}[1]{s_{#1}}     % Pure strategy


% ============================================================
\begin{document}

% Title
\begin{center}
  {\LARGE\bfseries\sffamily\color{defblue} MTH636 -- Game Theory}\\[6pt]
  {\Large\sffamily Lecture 21: Kuhn's Theorem \& SPE Existence}\\[4pt]
\end{center}
\vspace{0.5cm}
\hrule height 1.5pt
\vspace{0.3cm}

% ============================================================
\section{Existence of Nash Equilibria and SPE in Finite Games}

\begin{thmbox}{Kuhn's Theorem}
  Every extensive form game with \textbf{perfect information} has at least one Nash Equilibrium in pure strategies.
\end{thmbox}

\begin{thmbox}{SPE Existence Theorem}
  Every extensive form game with \textbf{perfect information} has at least one Subgame Perfect Equilibrium (SPE) in pure strategies.
\end{thmbox}

\begin{notebox}
  This SPE can be obtained using the \textbf{backward induction principle} -- i.e., the backward induction equilibrium is always a SPE.
\end{notebox}

% ============================================================
\section{Proof Sketch of SPE Existence via Backward Induction}

The proof proceeds by constructing a strategy profile $s^v$ via backward induction and then showing it constitutes a SPE. We outline the construction below.

\medskip
\noindent\textbf{Setup.} Consider the game tree rooted at $n_0$. At the root, Nature (or a player) chooses from branches $p_1, p_2, \ldots, p_k$, leading to subtrees rooted at $n_1, n_2, \ldots, n_k$, respectively. For each subtree rooted at $n_\ell$, we inductively obtain a strategy tuple:
\[
  s^{n_\ell} = (s_1^{n_\ell},\; \ldots,\; s_n^{n_\ell})
\]
where $s_i^{n_\ell}$ is player $i$'s strategy within the subtree rooted at $n_\ell$.

\subsection{Construction of the Equilibrium Strategy}

\textbf{Case (i): $n_0 \in V_0$ (Nature's node).}

For player $i$, define the strategy $s_i^v$ as follows: \emph{play $s_i^{n_\ell}$ if the game reaches $n_\ell$}. Let
\[
  s^v = (s_1^v,\; \ldots,\; s_n^v).
\]
We then show that $s^v$ is a SPE.

\medskip
\textbf{Case (ii): $n_0 \in V_{i_0}$, $i_0 \in N$ (some player's node).}

For player $i \neq i_0$: if the game reaches $n_\ell$, play $s_i^{n_\ell}$. Call this strategy $s_i^v$.

For player $i = i_0$: at node $n_0$, choose $n_{\ell_0}$ such that
\[
  U_{i_0}(s^{x_{\ell_0}}) \;\geq\; U_{i_0}(s^{n_\ell}) \quad \text{for all } \ell \in \{1, \ldots, k\},
\]
i.e., choose the branch $n_{\ell_0}$ that maximises player $i_0$'s payoff. Then if the game goes to $n_\ell$, follow strategy $s_{i_0}^{n_\ell}$.

\begin{notebox}
  The exact (formal) proof can be found in \textbf{Solani's book}.
\end{notebox}

% ============================================================
\section{Worked Example: SPE via Backward Induction}

Consider the following extensive form game. Nature (N) moves first at the root $n_0$, choosing Left ($L$) or Right ($R$). Player M (Man) then moves in each subtree.

\medskip
\noindent\textbf{Left subtree} (after $L$): Man (M) moves, choosing $F$ or $D$.
\begin{itemize}
  \item $F$: Woman (W) moves, choosing $F$ (payoff $(2,1)$) or $D$ (payoff $(0,0)$).
  \item $D$: Woman (W) moves, choosing $F$ (payoff $(0,0)$) or $D$ (payoff $(1,2)$).
\end{itemize}

\noindent\textbf{Right subtree} (after $R$): Woman (W) moves, choosing $F$ or $D$.
\begin{itemize}
  \item $F$: Man (M) moves, choosing $F$ (payoff $(2,1)$) or $D$ (payoff $(0,0)$).
  \item $D$: Man (M) moves, choosing $F$ (payoff $(0,0)$) or $D$ (payoff $(1,2)$).
\end{itemize}

\subsection{Backward Induction in the Left Subtree}

\noindent Applying backward induction in the left subtree (subgame rooted after $L$):
\begin{itemize}
  \item Woman's node after $M$ plays $F$: Woman chooses $F$ (payoff $1 > 0$). Resolved payoff: $(2,1)$.
  \item Woman's node after $M$ plays $D$: Woman chooses $D$ (payoff $2 > 0$). Resolved payoff: $(1,2)$.
  \item Man's node: compares $F \to (2,1)$ and $D \to (1,2)$. Man chooses $F$ (payoff $2 > 1$).
\end{itemize}
\noindent The backward induction equilibrium within the left subgame is $\bigl(F, (F,D)\bigr)$ with payoff $(2,1)$.

\subsection{Backward Induction in the Right Subtree}

\noindent Applying backward induction in the right subtree (subgame rooted after $R$):
\begin{itemize}
  \item Man's node after $W$ plays $F$: Man chooses $F$ (payoff $2 > 0$). Resolved payoff: $(2,1)$.
  \item Man's node after $W$ plays $D$: Man chooses $D$ (payoff $1 > 0$). Resolved payoff: $(1,2)$.
  \item Woman's node: compares $F \to (2,1)$ and $D \to (1,2)$. Woman chooses $D$ (payoff $2 > 1$).
\end{itemize}
\noindent The backward induction equilibrium within the right subgame is $\bigl(D, (F,D)\bigr)$ with payoff $(1,2)$.

\subsection{SPE for the Whole Game}

Combining the backward induction outcomes of both subgames:
\begin{itemize}
  \item Man's strategy: $M - (L,\; F,\; (F,D))$, i.e., play $F$ in the left subgame and $(F,D)$ in the right subgame.
  \item Woman's strategy: $W - (F,D,\; D)$, i.e., play $(F,D)$ in the left subgame and $D$ in the right subgame.
\end{itemize}

\noindent Concisely stated:
\[
  M:\; \bigl(L,\; F,\; (F,D)\bigr), \qquad W:\; (F,D,\; D)
\]
forms the \textbf{SPE for the whole game}.

\begin{notebox}
  Player M's expected utility: $U_M\bigl((L,F,(F,D)),\,(F,D,D)\bigr) = 2$.
\end{notebox}

% ============================================================
\section{Application: Chess and Determined Games}

\subsection{Structure of Chess as an Extensive Form Game}

Chess is a two-player, perfect information, zero-sum game. The game tree has:
\begin{itemize}
  \item A root node $O$.
  \item Subtrees for Player 1 ($P\text{-}1$), a ``Draw'' node ($D$), Player 2 ($P\text{-}2$), and Player 1 ($P\text{-}1$) again as the main branches.
\end{itemize}

The payoff matrix (from Player 1's perspective) takes the form:

\medskip
\begin{center}
\begin{tabular}{c|ccccc}
   & $s_2^1$ & $\cdots$ & $s_h^2$ \\\hline
  $s_1^1$ & $(1,-1)$ & $\cdots$ & \\
  $s_2^1$ & $(0,0)$ & $\cdots$ & \\
  $\vdots$ & & & \\
  $s_\ell^1$ & $(1,-1)$ & $(1_{\ell}=1)$ & $\cdots$ & $(1,-1)$ \\
  $\vdots$ & & & \\
  $s_m^1$ & $(-1,1)$ & $\cdots$ & \\
\end{tabular}
\end{center}

\noindent Here the payoffs represent: $(1,-1)$ = Player 1 wins, $(0,0)$ = Draw, $(-1,1)$ = Player 2 wins.

\subsection{Zermelo's Theorem Applied to Chess}

By the SPE existence theorem (Kuhn's Theorem), chess -- being a finite, two-player, perfect information game -- must have a pure strategy SPE. Because the game is zero-sum, this means:

\begin{thmbox}{Determinacy of Chess}
  Chess is a \textbf{determined} game: either Player 1 has a winning strategy, or Player 2 has a winning strategy, or both players can force a draw.
\end{thmbox}

Formally, the minimax values satisfy:
\[
  \underline{v} = \bar{v} = 1 \quad \text{(for Chess, the value exists)}
\]
where $\underline{v}$ is the maximin value and $\bar{v}$ is the minimax value. \textbf{For Chess}, we know:
\[
  \underline{v} = \bar{v} = 1
\]

\begin{notebox}
  While Zermelo's theorem guarantees that a winning (or drawing) strategy \emph{exists} for one of the players, it does \emph{not} tell us \emph{which} player has a winning strategy or what that strategy is. Finding such a strategy for Chess remains an open computational problem.
\end{notebox}

\begin{takebox}
  The key takeaway of Lectures 20--21 is the refinement chain:
  \[
    \text{SPE} \;\subset\; \text{Nash Equilibrium}
  \]
  SPE eliminates non-credible threats by requiring rationality at \emph{every} information set, not just along the equilibrium path. Backward induction provides a constructive method to find SPE in games with perfect information, and Kuhn's theorem guarantees such equilibria exist in pure strategies.
\end{takebox}

\end{document}
